\clearpage
\section{Candidate and Officeholder Characteristics}\label{sec:app_characteristics}
\suppressfloats[t]

Quota-seat winners are younger and have lower average education and reported assets than non-quota winners. The women-only Rajasthan comparisons show similar graduation rates and mean log assets across reservation categories. These comparisons describe the people selected under the two regimes, rather than measuring governing ability.

\subsection{Rajasthan measurement and samples}

The Rajasthan comparisons use the State Election Commission's 2020 candidate records linked to winners within the same GP and election phase. Candidate and event identifiers must be unique; winners additionally require a unique name match to the winner record. Age and occupation come from the candidate records. Graduation includes graduate, postgraduate, professional graduate, and professional postgraduate categories; missing education is excluded from that outcome. Reported unemployment is the recorded occupation category and need not imply an absence of work. We winsorize assets at the 10th and 90th percentiles of the candidate sample and take the natural logarithm of one plus assets.

For both states, ages outside 21--110 are treated as missing for age comparisons; other recorded characteristics remain usable. The lower bound follows the minimum age for office \citep{rajasthan1994act, up1947act}; the upper bound is a plausibility screen for recording errors. Table~\ref{tab:main_cand_char} compares all winners and candidates by reservation status. Table~\ref{tab:candidate_women} restricts each group to recorded women. This supplementary comparison helps interpret composition; it does not replace the comparison with all non-quota officeholders.

\begin{table}[!h]
\centering
\caption{\label{tab:candidate_women}Candidate Characteristics Among Women: Quota vs. Open Seats}
\centering
\fontsize{8}{10}\selectfont
\begin{threeparttable}
\begin{tabular}[t]{lllllll}
\toprule
\multicolumn{1}{c}{ } & \multicolumn{3}{c}{Women winners} & \multicolumn{3}{c}{Women candidates} \\
\cmidrule(l{3pt}r{3pt}){2-4} \cmidrule(l{3pt}r{3pt}){5-7}
Characteristic & Open & Quota & Difference & Open & Quota & Difference\\
\midrule
Age & 42.8 & 41.3 & 1.6 & 40.6 & 39.9 & 0.7\\
Graduate (\%) & 12.2 & 12.7 & -0.5 & 17.3 & 13.3 & 4.0\\
Unemployed (\%) & 53.1 & 61.1 & -8.0 & 54.8 & 61.6 & -6.8\\
Log assets & 12.6 & 12.6 & 0.0 & 12.1 & 12.1 & 0.0\\
\bottomrule
\end{tabular}
\begin{tablenotes}
\item \textit{Note: } 
\item Descriptive open-minus-quota comparisons, conditional on candidacy or winning. Graduation and unemployment are percentages; their differences are percentage points. Other rows use their stated units. Graduation rates exclude records with missing education. Age means exclude recorded ages outside 21--110 years; other outcomes retain those records. Assets winsorized at the 10th and 90th percentiles of the candidate sample before taking logs. These comparisons do not identify a causal effect on candidate quality.
\end{tablenotes}
\end{threeparttable}
\end{table}

\FloatBarrier

\clearpage
\subsection{Uttar Pradesh winner characteristics}

Table~\ref{tab:up_cand_char} uses the independent UP records described in Appendix~\ref{sec:weaver}. It compares winners' age and the source education measure in 2010, 2015, and 2021; assets are available only in the final wave and use the source's inverse-hyperbolic-sine transformation. Quota-seat winners have lower mean age and education in each election. These measures are reported in their source units and should not be compared directly with Rajasthan's graduation rate or log-assets measure.

\begin{table}
\centering
\caption{\label{tab:up_cand_char}Winner Characteristics: Quota vs. Open Seats (Uttar Pradesh)}
\centering
\fontsize{8}{10}\selectfont
\begin{threeparttable}
\begin{tabular}[t]{llllllllll}
\toprule
\multicolumn{1}{c}{ } & \multicolumn{3}{c}{2010} & \multicolumn{3}{c}{2015} & \multicolumn{3}{c}{2021} \\
\cmidrule(l{3pt}r{3pt}){2-4} \cmidrule(l{3pt}r{3pt}){5-7} \cmidrule(l{3pt}r{3pt}){8-10}
Characteristic & Open & Quota & Difference & Open & Quota & Difference & Open & Quota & Difference\\
\midrule
Age & 41.7 & 40.3 & 1.4 & 41.5 & 40.3 & 1.1 & 42.2 & 40.9 & 1.3\\
Education & 2.9 & 1.7 & 1.2 & 2.7 & 1.7 & 1.0 & 2.9 & 2.0 & 0.8\\
Assets (asinh) & -- & -- & -- & -- & -- & -- & 13.0 & 12.9 & 0.1\\
\bottomrule
\end{tabular}
\begin{tablenotes}
\item \textit{Note: } 
\item Descriptive comparisons using the Weaver data. Differences are open minus quota. The last wave is coded 2020 in the source labels and corresponds to the 2021 election. Assets are available only in that wave. Age means exclude recorded ages outside 21--110 years; other outcomes retain those records.
\end{tablenotes}
\end{threeparttable}
\end{table}

\FloatBarrier
